% TOP_PROBLEM: {"id": 1189, "title": "The Guralnick-Thompson Conjecture", "category_id": 4, "category": {"id": 4, "name": "algebra", "display_name": "Algebra", "description": "Group theory, ring theory, field theory, and algebraic structures.", "slug": "algebra", "order_index": 4, "created_at": "2026-07-31T15:26:25.671Z"}, "status": "open"}
% FIRST_POSTED: null
% ENTRY_KIND: statement_only
% Imported from: batch01.tex; UnsolvedMath ID 1189
\documentclass[11pt]{book}
\usepackage{fontspec}
\setmainfont{Latin Modern Roman}
\setsansfont{Latin Modern Sans}
\usepackage{amsmath,amssymb,mathtools}
\usepackage{unicode-math}
\setmathfont{Latin Modern Math}
\usepackage{xurl}
\usepackage[hidelinks]{hyperref}
\newcommand{\sref}[2]{\hyperlink{source:#1:#2}{[S#2]}}
\newcommand{\cataloguescope}[1]{\paragraph{Mathematical question.}#1}
\begin{document}
% UnsolvedMath ID 1189; source catalogue code ALG-025
\setcounter{section}{1188}
\section{The Guralnick-Thompson Conjecture}\label{1189}
{\small\sffamily UnsolvedMath ID 1189 · Algebra}
\subsection{Definitions and mathematical statement}

\paragraph{Shared notation.}$\mathbb N=\{1,2,\ldots\}$, and $\mathbb Z,\mathbb Q,\mathbb R,\mathbb C$ denote the integers, rationals, reals and complex numbers. Any other conventions are specified locally.


For $\sigma\in S_n$ let $\operatorname{ind}(\sigma)=n-c(\sigma)$, where $c(\sigma)$ counts all permutation cycles, including fixed points. A genus-zero system is a tuple $(\sigma_1,\ldots,\sigma_r)$ generating a transitive subgroup $G\leq S_n$ and satisfying
\[\sigma_1\cdots\sigma_r=1,\qquad\sum_{i=1}^r\operatorname{ind}(\sigma_i)=2(n-1).\]
Composition factors are the simple quotients in a composition series of $G$, considered up to group isomorphism. Let $E^*(0)$ be the union of these isomorphism types over all genus-zero systems, after excluding the alternating groups $A_m$ for $m\geq5$ and the cyclic groups $C_p$ of prime order.
\paragraph{Statement recorded in the supplied source.}
Classify the simple composition factors of all finite groups generated by genus-zero systems; equivalently, determine the exceptional set $E^*(0)$ in addition to the alternating and prime-order cyclic families. The named Guralnick--Thompson conjecture asserts finiteness of the analogous exceptional set at each fixed genus, defined by replacing $2(n-1)$ by $2(n+g-1)$; the cited source records that finiteness theorem as established. \sref{1189}{2}
\subsection{Short English statement}
Determine which finite simple groups occur as composition factors of genus-zero monodromy groups. These are groups generated by transitive permutation systems satisfying the genus-zero relation, with a historically proved finite-exception principle.
\subsection{Sources}
\begin{description}
\item[\hypertarget{source:1189:1}{S1}] ULAM.AI and the UnsolvedMath Contributors, \emph{UnsolvedMath: A Curated Collection of Open Mathematics Problems}, record 1189 (ALG-025). CC BY 4.0. \url{https://huggingface.co/datasets/ulamai/UnsolvedMath}.
\item[\hypertarget{source:1189:2}{S2}] Daniel Frohardt, Robert Guralnick and Kay Magaard, \emph{Primitive Monodromy Groups of Genus at most Two}, arXiv:1401.1303 (2014). Introduction and Definition1, pages1--2. \url{https://arxiv.org/pdf/1401.1303}.
\end{description}
\label{end:1189}
\end{document}
